We have demonstrated that simple game rules, multi-agent competition, and standard reinforcement learning algorithms at scale can induce agents to learn complex strategies and skills.
We observed emergence of as many as six distinct rounds of strategy and counter-strategy, suggesting that multi-agent self-play with simple game rules in sufficiently complex environments could lead to open-ended growth in complexity.
We then proposed to use transfer as a method to evaluate learning progress in open-ended environments and introduced a suite of targeted intelligence tests with which to compare agents in our domain.

Our results with hide-and-seek should be viewed as a proof of concept showing that multi-agent autocurricula can lead to physically grounded and human-relevant behavior.
We acknowledge that the strategy space in this environment is inherently bounded and likely will not surpass the six modes presented as is; however, because it is built in a high-fidelity physics simulator it is physically grounded and very extensible. In order to support further research in multi-agent autocurricula, we are open-sourcing our environment code.
% \footnote{Code can be found at \href{\codeurlclick}{\texttt{\codeurl}}.}

Hide-and-seek agents require an enormous amount of experience to progress through the six stages of emergence, likely because the reward functions are not directly aligned with the resulting behavior. While we have found that standard reinforcement learning algorithms are sufficient, reducing sample complexity in these systems will be an important line of future research. Better policy learning algorithms or policy architectures are orthogonal to our work and could be used to improve sample efficiency and performance on transfer evaluation metrics.

We also found that agents were very skilled at exploiting small inaccuracies in the design of the environment, such as seekers surfing on boxes without touching the ground, hiders running away from the environment while shielding themselves with boxes, or agents exploiting inaccuracies of the physics simulations to their advantage. Investigating methods to generate environments without these unwanted behaviors is another import direction of future research \citep{amodei2016, lehman2018}.